\nuevaReceta{Fideos con Sopa Miso}{2 personas}{25 min}{receta:sopa_miso}

    \begin{minipage}[t]{0.35\textwidth}
        \section*{Ingredientes}
        \begin{itemize}[leftmargin=*]
            \item 180 g de noodles o fideos de arroz
            \item 100 g de bacon o carne
            \item 2 cucharadas de pasta de miso
            \item 500 ml de agua
            \item 1 cucharada de salsa de soja
            \item 1/2 cebolla y 1 diente de ajo
            \item 80 g de guisantes o verdura verde
            \item 2 huevos cocidos
            \item \textbf{Toppings:} Cilantro, pimienta japonesa y sésamo
        \end{itemize}
    \end{minipage}
    \hfill
    \begin{minipage}[t]{0.6\textwidth}
        \section*{Preparación}
        \begin{enumerate}
            \item \textbf{Sofrito base:} En una olla con poco aceite, dorar la cebolla, el ajo y el bacon (o la carne elegida).
            \item \textbf{Reposo:} Cuando esté dorado, retirar la olla del fuego momentáneamente.
            \item \textbf{Caldo:} Calentar el agua y añadirla a la olla con los guisantes y la salsa de soja.
            \item \textbf{Cocción:} Llevar a ebullición y añadir los fideos hasta que estén en su punto. Retirar del fuego, disolver aparte el miso con un poco de caldo e incorporarlo sin volver a hervir.
            \item \textbf{Emplatado:} Servir en un bol hondo y decorar con cilantro fresco, pimienta japonesa, semillas de sésamo y el huevo cocido.
        \end{enumerate}
    \end{minipage}

    \vspace{1em}
    \begin{tcolorbox}[colback=orange!5, colframe=orange!50, title=Nota, size=small]
        Para un huevo cocido perfecto (tipo ramen), déjalo hervir exactamente 6-7 minutos y pásalo por agua con hielo.
    \end{tcolorbox}
\end{receta}
